\documentclass[11pt,a4paper]{article}

\usepackage[utf8]{inputenc}
\usepackage{amsmath,amssymb,amsfonts}
\usepackage{geometry}
\geometry{left=2.5cm, right=2.5cm, top=2.8cm, bottom=2.8cm}
\usepackage{xcolor}
\usepackage{booktabs}
\usepackage{fancyhdr}
\usepackage{listings}
\usepackage{hyperref}
\usepackage{enumitem}

% Paleta institucional UAN
\definecolor{uanblue}{RGB}{0, 45, 98}
\definecolor{uancyan}{RGB}{0, 140, 200}
\definecolor{accentgreen}{RGB}{0, 160, 120}
\definecolor{codebg}{RGB}{248, 249, 252}
\definecolor{bordercolor}{RGB}{210, 220, 235}
\definecolor{darkgray}{RGB}{60, 64, 70}

\hypersetup{
    colorlinks=true,
    linkcolor=uanblue,
    citecolor=uancyan,
    urlcolor=uancyan,
    pdftitle={Aurora By Juan Esteban Vera - Anexo de Codigo Fuente y Arquitectura DSP},
    pdfauthor={Juan Esteban Vera Rodriguez}
}

\pagestyle{fancy}
\fancyhf{}
\fancyhead[L]{\small \textcolor{uanblue}{\textbf{Universidad Antonio Nariño}} · Sede Ibagué · Ingeniería Electrónica}
\fancyhead[R]{\small \textcolor{gray}{Aurora By Juan Esteban Vera · Código Fuente DSP}}
\fancyfoot[L]{\small \textcolor{gray}{Anexo Técnico 02 — Arquitectura de Software y Código Fuente}}
\fancyfoot[R]{\small \thepage}
\renewcommand{\headrulewidth}{0.6pt}
\renewcommand{\footrulewidth}{0.4pt}

\lstset{
    language=Python,
    backgroundcolor=\color{codebg},
    basicstyle=\ttfamily\scriptsize,
    breaklines=true,
    captionpos=b,
    frame=single,
    rulecolor=\color{bordercolor},
    keywordstyle=\color{uanblue}\bfseries,
    commentstyle=\color{gray}\itshape,
    stringstyle=\color{accentgreen},
    showstringspaces=false,
    tabsize=4
}

\begin{document}

% =========================================================================
% PORTADA INSTITUCIONAL
% =========================================================================
\begin{titlepage}
    \centering
    \vspace*{0.5cm}
    
    {\scshape\LARGE \textbf{UNIVERSIDAD ANTONIO NARIÑO}}\\[0.3cm]
    {\scshape\large FACULTAD DE INGENIERÍA · PROGRAMA DE INGENIERÍA ELECTRÓNICA}\\[0.15cm]
    {\small SEDE IBAGUÉ — CAPITAL MUSICAL DE COLOMBIA}\\[1.5cm]
    
    \rule{\linewidth}{1.8pt}\\[0.4cm]
    {\Huge \bfseries \textcolor{uanblue}{ANEXO DE CÓDIGO FUENTE}}\\[0.35cm]
    {\Large \bfseries ARQUITECTURA DE SOFTWARE DSP, ALGORITMOS DE SÍNTESIS Y MOTOR DE AUDIO EN TIEMPO REAL}\\[0.25cm]
    {\large \itshape Aurora By Juan Esteban Vera — Yamaha MONTAGE DSP Edition}\\[0.4cm]
    \rule{\linewidth}{1.8pt}\\[2.0cm]
    
    \begin{minipage}{0.85\textwidth}
        \large
        \textbf{Autor y Desarrollador:}\\[0.15cm]
        \textbf{Juan Esteban Vera Rodríguez}\\
        Investigador Principal · Estudiante de Ingeniería Electrónica\\
        Universidad Antonio Nariño · Sede Ibagué\\[0.8cm]
        
        \textbf{Proyecto Institucional:}\\
        Repositorio de Código Fuente y Registro de Propiedad Intelectual\\
        Dirección Nacional de Derecho de Autor (DNDA)\\[0.8cm]
        
        \textbf{Identificador Documental:}\\
        ANX-SRC-002 · Especificación de Software y Códigos de Ingeniería
    \end{minipage}
    
    \vfill
    {\large Ibagué, Tolima, Colombia · 2026}
\end{titlepage}

\newpage
\tableofcontents
\newpage

% =========================================================================
\section{Arquitectura Global del Software}
% =========================================================================

El sistema \textbf{Aurora By Juan Esteban Vera} fue diseñado bajo un patrón arquitectónico desacoplado y orientado a objetos, garantizando que el cómputo matemático de audio en tiempo real opere de forma estrictamente asíncrona e independiente del hilo principal de renderizado gráfico (GUI).

\subsection{Diagrama de Componentes y Flujo de Señales}
\begin{itemize}[leftmargin=1.5cm]
    \item \textbf{Capa de Síntesis y DSP (\texttt{synth/engine.py}):}
    Aloja las estructuras inmutables de parámetros (\texttt{SynthParameters}), el generador de formas de onda con PolyBLEP, la modulación FM-X, el filtro TPT SVF, la envolvente ADSR vectorizada, y la cadena de efectos (Drive, Stereo Delay, Reverb Schroeder).
    \item \textbf{Capa de Procesamiento Vocal (\texttt{synth/vocal.py}):}
    Contiene el detector de tono fundamental (FFT / autocorrelación normalizada) y el corrector granular de tono afinado en tiempo real según escalas modales.
    \item \textbf{Capa de Presets de Fábrica (\texttt{synth/presets.py}):}
    Define los bancos de memoria tímbrica calibrados con base en los instrumentos clásicos de Yamaha (DX7, CS-80 y MONTAGE).
    \item \textbf{Capa de Interfaz de Usuario y Secuenciador (\texttt{synth/app.py}):}
    Implementa la consola de control con el teclado extendido de 25 teclas, el dial macro Super Knob, el secuenciador de 16 pasos, el osciloscopio en vivo y los medidores balísticos.
    \item \textbf{Capa de Validación Automatizada (\texttt{tests/test\_engine.py}):}
    Suite de 13 pruebas unitarias para verificación de estabilidad numérica y latencia.
\end{itemize}

% =========================================================================
\section{Módulo 1: Motor DSP en Tiempo Real (\texttt{synth/engine.py})}
% =========================================================================

\subsection{Estructura Inmutable de Parámetros}
La clase \texttt{SynthParameters} encapsula los 17 parámetros de control tímbrico y espacial mediante el decorador \texttt{@dataclass(frozen=True)}, asegurando concurrencia \textit{thread-safe} sin bloqueos (\textit{lock-free}):

\begin{lstlisting}
@dataclass(frozen=True)
class SynthParameters:
    # Osciladores principales
    waveform_1: str = "sine"
    waveform_2: str = "triangle"
    osc2_mix: float = 0.18
    detune_cents: float = 0.0
    osc2_semitones: float = 0.0

    # Sintesis FM (Yamaha DX / FM-X style)
    fm_depth: float = 0.0

    # Envolvente ADSR
    attack: float = 0.012
    decay: float = 0.18
    sustain: float = 0.72
    release: float = 0.32

    # Filtro & Resonancia (Yamaha CS-80 style)
    cutoff: float = 6500.0
    resonance: float = 0.0

    # Procesador de Efectos (Yamaha MONTAGE Insertion FX)
    drive: float = 0.0
    delay_time: float = 0.25
    delay_feedback: float = 0.35
    delay_mix: float = 0.0

    # Modulacion, espacio y volumen
    vibrato_rate: float = 5.0
    vibrato_depth: float = 0.0
    master_gain: float = 0.35
    reverb_mix: float = 0.18
    reverb_room: float = 0.55

    # Mando maestro Motion Control (Yamaha MONTAGE Super Knob)
    super_knob: float = 0.0
\end{lstlisting}

\subsection{Ecuación de Afinación y Supresión de Aliasing PolyBLEP}
La afinación temperada igual sigue la norma ISO 16, mientras que el algoritmo PolyBLEP inserta correcciones polinómicas para suavizar las discontinuidades en los bordes de fase:

\begin{lstlisting}
def midi_to_hz(note: int | float) -> float:
    return 440.0 * 2.0 ** ((note - 69) / 12.0)

def _poly_blep(phase: np.ndarray, increment: float) -> np.ndarray:
    result = np.zeros_like(phase)
    first = phase < increment
    t = phase[first] / increment
    result[first] = 2.0 * t - t * t - 1.0
    last = phase > 1.0 - increment
    t = (phase[last] - 1.0) / increment
    result[last] = t * t + 2.0 * t + 1.0
    return result
\end{lstlisting}

\subsection{Generador de Formas de Onda y Modulación FM-X}
El renderizador vectorial genera formas de onda sinusoidales, triangulares, cuadradas y diente de sierra, incorporando modulación de fase para la síntesis FM:

\begin{lstlisting}
def _render_waveform(waveform: str, phase: np.ndarray, increment: float, phase_mod: np.ndarray | None = None) -> np.ndarray:
    total_phase = phase if phase_mod is None else (phase + phase_mod)
    total_phase = np.mod(total_phase, 1.0)

    if waveform == "sine":
        return np.sin(2.0 * np.pi * total_phase)
    if waveform == "triangle":
        saw = 2.0 * total_phase - 1.0
        return 2.0 * np.abs(saw) - 1.0
    if waveform == "saw":
        naive = 2.0 * total_phase - 1.0
        return naive - _poly_blep(total_phase, increment)
    if waveform == "square":
        naive = np.where(total_phase < 0.5, 1.0, -1.0)
        corr_0 = _poly_blep(total_phase, increment)
        corr_half = _poly_blep(np.mod(total_phase + 0.5, 1.0), increment)
        return naive + corr_0 - corr_half
    if waveform == "noise":
        return np.random.uniform(-1.0, 1.0, size=phase.shape)
    return np.sin(2.0 * np.pi * total_phase)
\end{lstlisting}

\subsection{Filtro Analógico Virtual TPT SVF de 2 Polos}
Filtro de variable de estado discretizado mediante Transformada Bilineal con prewarping de frecuencia:

\begin{lstlisting}
def _process_svf(signal: np.ndarray, cutoff_hz: float, resonance: float, sample_rate: int, s1: float, s2: float):
    cutoff = np.clip(cutoff_hz, 40.0, sample_rate * 0.45)
    res = np.clip(resonance, 0.0, 0.95)
    g = math.tan(math.pi * cutoff / sample_rate)
    k = 2.0 * (1.0 - res)
    a1 = 1.0 / (1.0 + g * (g + k))
    a2 = g * a1
    a3 = g * a2

    out = np.empty_like(signal)
    for n in range(signal.shape[0]):
        v0 = signal[n]
        v1 = a1 * s1 + a2 * (v0 - s2)
        v2 = s2 + a3 * (v0 - s2) + a2 * s1
        s1 = 2.0 * v1 - s1
        s2 = 2.0 * v2 - s2
        out[n] = math.tanh(v2)
    return out, s1, s2
\end{lstlisting}

\subsection{Ciclo de Vida de las Voces y Envolvente ADSR Vectorizada}
Cada voz activa calcula su ganancia temporal en función de las 4 etapas ADSR:

\begin{lstlisting}
class ActiveVoice:
    def __init__(self, note: int, sample_rate: int, params: SynthParameters):
        self.note = note
        self.sample_rate = sample_rate
        self.params = params
        self.base_frequency = midi_to_hz(note)
        self.phase_1 = 0.0
        self.phase_2 = 0.0
        self.lfo_phase = 0.0
        self.filter_s1_left = 0.0
        self.filter_s2_left = 0.0
        self.filter_s1_right = 0.0
        self.filter_s2_right = 0.0
        self.age_samples = 0
        self.released = False
        self.release_start_sample = 0
        self.release_start_level = 1.0
        self.is_finished = False

    def render(self, num_samples: int, params: SynthParameters) -> np.ndarray:
        # Generacion de incrementos angulares
        dt1 = self.base_frequency / self.sample_rate
        osc2_ratio = 2.0 ** (params.osc2_semitones / 12.0 + params.detune_cents / 1200.0)
        dt2 = (self.base_frequency * osc2_ratio) / self.sample_rate

        phase_1_arr = np.mod(self.phase_1 + np.arange(num_samples) * dt1, 1.0)
        phase_2_arr = np.mod(self.phase_2 + np.arange(num_samples) * dt2, 1.0)

        # Modulador FM
        modulator = _render_waveform(params.waveform_2, phase_2_arr, dt2)
        fm_term = modulator * (params.fm_depth * 1.8) if params.fm_depth > 0.001 else None
        carrier = _render_waveform(params.waveform_1, phase_1_arr, dt1, fm_term)

        # Mezcla balanceada
        raw_voice = carrier * (1.0 - params.osc2_mix * 0.5) + modulator * (params.osc2_mix * 0.7)

        # Calculo de Envolvente ADSR
        env = self._calculate_adsr(num_samples, params)
        voice_out = raw_voice * env

        # Filtrado SVF
        filtered, self.filter_s1_left, self.filter_s2_left = _process_svf(
            voice_out, params.cutoff, params.resonance, self.sample_rate,
            self.filter_s1_left, self.filter_s2_left
        )

        # Avance de acumuladores
        self.phase_1 = (self.phase_1 + num_samples * dt1) % 1.0
        self.phase_2 = (self.phase_2 + num_samples * dt2) % 1.0
        self.age_samples += num_samples

        stereo = np.column_stack((filtered, filtered))
        return stereo
\end{lstlisting}

\subsection{El Motion Control Super Knob y Efectos Master}
La función \texttt{\_effective\_params} calcula la interpolación macro entre los parámetros base y el dial Super Knob:

\begin{lstlisting}
def _effective_params(params: SynthParameters) -> SynthParameters:
    knob = float(np.clip(params.super_knob, 0.0, 1.0))
    if knob <= 0.0001:
        return params

    return replace(
        params,
        cutoff=float(np.clip(params.cutoff * (1.0 + 2.0 * knob), 60.0, 15000.0)),
        resonance=float(np.clip(params.resonance + 0.35 * knob, 0.0, 0.94)),
        fm_depth=float(np.clip(params.fm_depth + 0.55 * knob, 0.0, 1.0)),
        drive=float(np.clip(params.drive + 0.45 * knob, 0.0, 1.0)),
        delay_mix=float(np.clip(params.delay_mix + 0.30 * knob, 0.0, 0.85)),
        reverb_mix=float(np.clip(params.reverb_mix + 0.35 * knob, 0.0, 0.85)),
    )
\end{lstlisting}

% =========================================================================
\section{Módulo 2: Procesador Vocal \& Autotune (\texttt{synth/vocal.py})}
% =========================================================================

\subsection{Detección de Tono Fundamental (Pitch Detector)}
El algoritmo combina filtrado paso-bajo y autocorrelación normalizada mediante FFT (Teorema de Wiener-Khinchin):

\begin{lstlisting}
class PitchDetector:
    def __init__(self, sample_rate: int = 44100, buffer_size: int = 4096):
        self.sample_rate = sample_rate
        self.buffer_size = buffer_size

    def detect_pitch(self, audio_buffer: np.ndarray) -> tuple[float, float]:
        if len(audio_buffer) < self.buffer_size:
            return 0.0, 0.0

        # Normalizacion y ventana Hann
        signal = audio_buffer[-self.buffer_size:]
        signal = signal - np.mean(signal)
        rms = np.sqrt(np.mean(signal**2))
        if rms < 0.008:
            return 0.0, 0.0

        windowed = signal * np.hanning(len(signal))
        # Autocorrelacion via FFT
        fft_data = np.fft.rfft(windowed, n=self.buffer_size * 2)
        corr = np.fft.irfft(fft_data * np.conj(fft_data))[:self.buffer_size]
        corr = corr / (corr[0] + 1e-12)

        # Rango de busqueda: 70 Hz a 1000 Hz
        min_lag = int(self.sample_rate / 1000.0)
        max_lag = int(self.sample_rate / 70.0)

        peak_lag = min_lag + np.argmax(corr[min_lag:max_lag])
        confidence = float(corr[peak_lag])

        if confidence < 0.35:
            return 0.0, confidence

        detected_hz = float(self.sample_rate / peak_lag)
        return detected_hz, confidence
\end{lstlisting}

\subsection{Cuantización de Escala y Corrector Granular}
\begin{lstlisting}
SCALE_PRESETS = {
    "C Mayor": [0, 2, 4, 5, 7, 9, 11],
    "C Menor": [0, 2, 3, 5, 7, 8, 10],
    "A Menor": [0, 2, 3, 5, 7, 8, 10],
    "C Cromatica": list(range(12)),
}

def snap_to_scale(midi_pitch: float, scale_intervals: list[int], root: int = 0) -> float:
    semitone = round(midi_pitch)
    note_class = (semitone - root) % 12
    closest = min(scale_intervals, key=lambda p: min(abs(note_class - p), 12 - abs(note_class - p)))
    octave = (semitone - root) // 12
    return float(root + octave * 12 + closest)
\end{lstlisting}

% =========================================================================
\section{Módulo 3: Banco de Presets Yamaha (\texttt{synth/presets.py})}
% =========================================================================

\begin{lstlisting}
YAMAHA_PRESETS: dict[str, SynthParameters] = {
    "Yamaha DX7 E-Piano": SynthParameters(
        waveform_1="sine", waveform_2="sine", osc2_mix=0.45,
        osc2_semitones=12.0, detune_cents=2.0, fm_depth=0.42,
        attack=0.005, decay=0.65, sustain=0.18, release=0.45,
        cutoff=8200.0, resonance=0.05, vibrato_rate=5.5,
        vibrato_depth=0.04, drive=0.08, delay_mix=0.22,
        reverb_mix=0.32, master_gain=0.38
    ),
    "CS-80 Poly Brass": SynthParameters(
        waveform_1="saw", waveform_2="saw", osc2_mix=0.55,
        osc2_semitones=0.0, detune_cents=12.0, fm_depth=0.0,
        attack=0.085, decay=0.45, sustain=0.75, release=0.55,
        cutoff=3400.0, resonance=0.42, vibrato_rate=4.8,
        vibrato_depth=0.14, drive=0.12, reverb_mix=0.40,
        master_gain=0.34
    ),
    "Montage Super Lead": SynthParameters(
        waveform_1="saw", waveform_2="square", osc2_mix=0.60,
        osc2_semitones=7.0, detune_cents=14.0, fm_depth=0.35,
        attack=0.008, decay=0.25, sustain=0.85, release=0.35,
        cutoff=7800.0, resonance=0.38, drive=0.35,
        delay_mix=0.35, reverb_mix=0.28, super_knob=0.35,
        master_gain=0.32
    ),
    "Ambient Space Pad": SynthParameters(
        waveform_1="triangle", waveform_2="saw", osc2_mix=0.40,
        osc2_semitones=12.0, detune_cents=8.0, fm_depth=0.12,
        attack=0.45, decay=0.85, sustain=0.80, release=1.20,
        cutoff=2800.0, resonance=0.25, delay_mix=0.45,
        reverb_mix=0.65, reverb_room=0.85, master_gain=0.30
    ),
    "Analog Bass 80s": SynthParameters(
        waveform_1="saw", waveform_2="square", osc2_mix=0.50,
        osc2_semitones=-12.0, detune_cents=5.0, fm_depth=0.05,
        attack=0.006, decay=0.22, sustain=0.40, release=0.18,
        cutoff=2200.0, resonance=0.52, drive=0.28,
        delay_mix=0.0, reverb_mix=0.10, master_gain=0.42
    ),
    "Square Bell FM": SynthParameters(
        waveform_1="sine", waveform_2="square", osc2_mix=0.40,
        osc2_semitones=19.0, detune_cents=0.0, fm_depth=0.68,
        attack=0.002, decay=0.85, sustain=0.05, release=0.60,
        cutoff=9000.0, resonance=0.10, delay_mix=0.30,
        reverb_mix=0.45, master_gain=0.32
    ),
    "Vocal AutoTune FX": SynthParameters(
        waveform_1="saw", waveform_2="sine", osc2_mix=0.30,
        osc2_semitones=0.0, detune_cents=4.0, fm_depth=0.10,
        attack=0.010, decay=0.30, sustain=0.70, release=0.25,
        cutoff=6000.0, resonance=0.20, drive=0.15,
        delay_mix=0.25, reverb_mix=0.30, master_gain=0.35
    )
}
\end{lstlisting}

% =========================================================================
\section{Módulo 4: Consola de Trabajo GUI (\texttt{synth/app.py})}
% =========================================================================

\subsection{El Componente Gráfico Super Knob}
Widget personalizado renderizado mediante primitivas vectoriales de Tkinter Canvas con soporte para arrastre y rueda del mouse:

\begin{lstlisting}
class SuperKnob(tk.Canvas):
    def __init__(self, master, size=150, command=None, **kwargs):
        super().__init__(master, width=size, height=size, bg="#0d141e", highlightthickness=0, **kwargs)
        self.size = size
        self.command = command
        self.value = 0.0
        self._dragging = False
        self._last_y = 0

        self.bind("<Button-1>", self._on_press)
        self.bind("<B1-Motion>", self._on_drag)
        self.bind("<ButtonRelease-1>", self._on_release)
        self.bind("<MouseWheel>", self._on_wheel)
        self._draw()

    def set_value(self, val: float, notify: bool = True):
        self.value = float(np.clip(val, 0.0, 1.0))
        self._draw()
        if notify and self.command:
            self.command(self.value)

    def _draw(self):
        self.delete("all")
        cx = cy = self.size / 2.0
        r_outer = self.size * 0.42
        r_inner = self.size * 0.32

        # Anillo de fondo y arco de progreso activo
        self.create_oval(cx - r_outer, cy - r_outer, cx + r_outer, cy + r_outer, outline="#1a2738", width=7)
        start_deg = -225
        extent_deg = -270 * self.value
        self.create_arc(cx - r_outer, cy - r_outer, cx + r_outer, cy + r_outer,
                        start=start_deg, extent=extent_deg, outline="#00e5ff", width=7, style="arc")

        # Dial central y puntero indicador
        self.create_oval(cx - r_inner, cy - r_inner, cx + r_inner, cy + r_inner, fill="#131e2d", outline="#253549", width=2)
        angle_rad = math.radians(225 - 270 * self.value)
        px = cx + (r_inner * 0.75) * math.cos(angle_rad)
        py = cy - (r_inner * 0.75) * math.sin(angle_rad)
        self.create_line(cx, cy, px, py, fill="#00e5ff", width=3, capstyle=tk.ROUND)

        # Etiqueta de porcentaje
        self.create_text(cx, cy, text=f"{int(self.value * 100)}%", fill="#ffffff", font=("Segoe UI", 12, "bold"))
\end{lstlisting}

\subsection{Teclado Extendido de 25 Teclas Interactivas}
Implementación de doble octava continua con retroalimentación visual al tacto:

\begin{lstlisting}
KEY_SPECS_25 = [
    # Octava 1 (C3 - B3)
    {"midi_offset": 0, "name": "C", "is_black": False, "key": "Z"},
    {"midi_offset": 1, "name": "C#", "is_black": True,  "key": "S"},
    {"midi_offset": 2, "name": "D", "is_black": False, "key": "X"},
    {"midi_offset": 3, "name": "D#", "is_black": True,  "key": "D"},
    {"midi_offset": 4, "name": "E", "is_black": False, "key": "C"},
    {"midi_offset": 5, "name": "F", "is_black": False, "key": "V"},
    {"midi_offset": 6, "name": "F#", "is_black": True,  "key": "G"},
    {"midi_offset": 7, "name": "G", "is_black": False, "key": "B"},
    {"midi_offset": 8, "name": "G#", "is_black": True,  "key": "H"},
    {"midi_offset": 9, "name": "A", "is_black": False, "key": "N"},
    {"midi_offset": 10, "name": "A#", "is_black": True, "key": "J"},
    {"midi_offset": 11, "name": "B", "is_black": False, "key": "M"},
    # Octava 2 (C4 - C5)
    {"midi_offset": 12, "name": "C", "is_black": False, "key": "Q"},
    {"midi_offset": 13, "name": "C#", "is_black": True, "key": "2"},
    {"midi_offset": 14, "name": "D", "is_black": False, "key": "W"},
    {"midi_offset": 15, "name": "D#", "is_black": True, "key": "3"},
    {"midi_offset": 16, "name": "E", "is_black": False, "key": "E"},
    {"midi_offset": 17, "name": "F", "is_black": False, "key": "R"},
    {"midi_offset": 18, "name": "F#", "is_black": True, "key": "5"},
    {"midi_offset": 19, "name": "G", "is_black": False, "key": "T"},
    {"midi_offset": 20, "name": "G#", "is_black": True, "key": "6"},
    {"midi_offset": 21, "name": "A", "is_black": False, "key": "Y"},
    {"midi_offset": 22, "name": "A#", "is_black": True, "key": "7"},
    {"midi_offset": 23, "name": "B", "is_black": False, "key": "U"},
    {"midi_offset": 24, "name": "C", "is_black": False, "key": "I"},
]
\end{lstlisting}

% =========================================================================
\section{Módulo 5: Suite de Pruebas Unitarias (\texttt{tests/test\_engine.py})}
% =========================================================================

\begin{lstlisting}
class TestAuroraEngine(unittest.TestCase):
    def setUp(self):
        self.engine = AudioEngine(sample_rate=44100, buffer_size=256)

    def tearDown(self):
        self.engine.stop()

    def test_a4_frequency(self):
        self.assertAlmostEqual(midi_to_hz(69), 440.0, places=4)

    def test_note_generates_stereo_audio(self):
        self.engine.note_on(60)
        chunk = self.engine.process_block(256)
        self.assertEqual(chunk.shape, (256, 2))
        self.assertTrue(np.max(np.abs(chunk)) > 0.0)

    def test_release_finishes_voice(self):
        self.engine.note_on(60)
        self.engine.note_off(60)
        # Avanzar el tiempo suficiente para completar el release
        for _ in range(30):
            self.engine.process_block(256)
        self.assertEqual(len(self.engine.voices), 0)

    def test_pitch_detector_reads_a4(self):
        detector = PitchDetector(sample_rate=44100, buffer_size=4096)
        t = np.arange(4096) / 44100.0
        sine_440 = 0.8 * np.sin(2.0 * np.pi * 440.0 * t)
        freq, conf = detector.detect_pitch(sine_440)
        self.assertAlmostEqual(freq, 440.0, delta=3.0)
        self.assertGreater(conf, 0.25)

    def test_scale_snaps_to_allowed_note(self):
        c_major = SCALE_PRESETS["C Mayor"]
        snapped = snap_to_scale(61.0, c_major) # C# -> C
        self.assertEqual(snapped, 60.0)

    def test_resonant_filter_stability(self):
        noise = np.random.uniform(-1.0, 1.0, 512)
        out, s1, s2 = _process_svf(noise, cutoff_hz=1000.0, resonance=0.92,
                                  sample_rate=44100, s1=0.0, s2=0.0)
        self.assertFalse(np.any(np.isnan(out)))
        self.assertFalse(np.any(np.isinf(out)))
\end{lstlisting}

% =========================================================================
\section{Punto de Entrada Principal (\texttt{main.py})}
% =========================================================================

\begin{lstlisting}
from synth.app import main

if __name__ == "__main__":
    main()
\end{lstlisting}

\end{document}
